% Options for packages loaded elsewhere
\PassOptionsToPackage{unicode}{hyperref}
\PassOptionsToPackage{hyphens}{url}
\PassOptionsToPackage{dvipsnames,svgnames*,x11names*}{xcolor}
%
\documentclass[
  10pt,
]{article}
\usepackage{lmodern}
\usepackage{amssymb,amsmath}
\usepackage{ifxetex,ifluatex}
\ifnum 0\ifxetex 1\fi\ifluatex 1\fi=0 % if pdftex
  \usepackage[T1]{fontenc}
  \usepackage[utf8]{inputenc}
  \usepackage{textcomp} % provide euro and other symbols
\else % if luatex or xetex
  \usepackage{unicode-math}
  \defaultfontfeatures{Scale=MatchLowercase}
  \defaultfontfeatures[\rmfamily]{Ligatures=TeX,Scale=1}
  \setmainfont[]{DejaVu Serif}
  \setmonofont[]{DejaVu Sans Mono}
\fi
% Use upquote if available, for straight quotes in verbatim environments
\IfFileExists{upquote.sty}{\usepackage{upquote}}{}
\IfFileExists{microtype.sty}{% use microtype if available
  \usepackage[]{microtype}
  \UseMicrotypeSet[protrusion]{basicmath} % disable protrusion for tt fonts
}{}
\makeatletter
\@ifundefined{KOMAClassName}{% if non-KOMA class
  \IfFileExists{parskip.sty}{%
    \usepackage{parskip}
  }{% else
    \setlength{\parindent}{0pt}
    \setlength{\parskip}{6pt plus 2pt minus 1pt}}
}{% if KOMA class
  \KOMAoptions{parskip=half}}
\makeatother
\usepackage{xcolor}
\IfFileExists{xurl.sty}{\usepackage{xurl}}{} % add URL line breaks if available
\IfFileExists{bookmark.sty}{\usepackage{bookmark}}{\usepackage{hyperref}}
\hypersetup{
  colorlinks=true,
  linkcolor=Maroon,
  filecolor=Maroon,
  citecolor=Blue,
  urlcolor=Blue,
  pdfcreator={LaTeX via pandoc}}
\urlstyle{same} % disable monospaced font for URLs
\usepackage[margin=0.6in,landscape]{geometry}
\usepackage{longtable,booktabs}
% Correct order of tables after \paragraph or \subparagraph
\usepackage{etoolbox}
\makeatletter
\patchcmd\longtable{\par}{\if@noskipsec\mbox{}\fi\par}{}{}
\makeatother
% Allow footnotes in longtable head/foot
\IfFileExists{footnotehyper.sty}{\usepackage{footnotehyper}}{\usepackage{footnote}}
\makesavenoteenv{longtable}
\setlength{\emergencystretch}{3em} % prevent overfull lines
\providecommand{\tightlist}{%
  \setlength{\itemsep}{0pt}\setlength{\parskip}{0pt}}
\setcounter{secnumdepth}{-\maxdimen} % remove section numbering
\renewcommand{\arraystretch}{1.6}

\author{}
\date{}

\begin{document}

\hypertarget{phase-0.0-slide-rule-basics}{%
\section{Phase 0.0 --- Slide Rule
Basics}\label{phase-0.0-slide-rule-basics}}

\textbf{Type:} Tool drill, no mission narrative. \textbf{Target time:}
\textasciitilde30 min if the scales are already comfortable; up to an
afternoon if not. \textbf{Prerequisite:} None. This is the dependency
root for everything else in the campaign. \textbf{Assumes:} A generic
linear slide rule with C, D, A, B, K, S, and T scales (a standard 10"
engineering rule --- Pickett, K\&E, Post, etc. all work). This is the
baseline for Phase 0; nothing stops a later stage from calling for a
specialized rule (vector rule, log-log duplex for exponentials, a period
Trig/Kepler-specific rule, etc.) if the era or problem genuinely
warrants it --- just note the swap explicitly in that problem's file
when it happens.

\hypertarget{a-note-before-you-start}{%
\subsection{A note before you start}\label{a-note-before-you-start}}

You've got E6B experience, not linear-rule experience --- treat that as
a \textbf{false friend}, not a head start. The E6B is a \emph{circular}
rule built around wind-triangle and fuel-burn conventions (rotating
compass rose, specialized scales for TAS/groundspeed). A linear C/D/A/B
rule shares the underlying logarithmic-scale principle but none of the
E6B's specialized layout. Two habits from the E6B will actively work
against you here:

\begin{itemize}
\tightlist
\item
  \textbf{No fixed ``index'' story.} On the E6B you're often reading off
  a wind triangle with a consistent visual frame. On a linear rule, the
  C and D scales slide relative to each other and there's no rotating
  dial to anchor your eye --- you have to track which index (1 at the
  left end of C, or the 10 at the right end) you're using on a given
  problem, and that choice changes every time.
\item
  \textbf{Decimal placement is entirely on you.} The E6B's problems are
  bounded (airspeeds, fuel flows) so wrong-order-of-magnitude answers
  are obviously wrong at a glance. A slide rule reads the same for 2.35,
  23.5, and 235 --- the scale only gives you the \emph{significant
  digits}. You have to track the decimal point yourself with a quick
  mental estimate before you trust the rule's answer. This is arguably
  the single most important habit in this whole drill, and it's the one
  place hand computation fails silently if you skip it.
\end{itemize}

\hypertarget{scale-reference}{%
\subsection{Scale reference}\label{scale-reference}}

\begin{longtable}[]{@{}lll@{}}
\toprule
\begin{minipage}[b]{0.28\columnwidth}\raggedright
Scale\strut
\end{minipage} & \begin{minipage}[b]{0.44\columnwidth}\raggedright
What it is\strut
\end{minipage} & \begin{minipage}[b]{0.20\columnwidth}\raggedright
Use\strut
\end{minipage}\tabularnewline
\midrule
\endhead
\begin{minipage}[t]{0.28\columnwidth}\raggedright
C, D\strut
\end{minipage} & \begin{minipage}[t]{0.44\columnwidth}\raggedright
Identical logarithmic scales, 1--10\strut
\end{minipage} & \begin{minipage}[t]{0.20\columnwidth}\raggedright
Multiplication, division\strut
\end{minipage}\tabularnewline
\begin{minipage}[t]{0.28\columnwidth}\raggedright
A, B\strut
\end{minipage} & \begin{minipage}[t]{0.44\columnwidth}\raggedright
Two decades (1--100) squeezed into the C/D length\strut
\end{minipage} & \begin{minipage}[t]{0.20\columnwidth}\raggedright
Squares, square roots\strut
\end{minipage}\tabularnewline
\begin{minipage}[t]{0.28\columnwidth}\raggedright
K\strut
\end{minipage} & \begin{minipage}[t]{0.44\columnwidth}\raggedright
Three decades (1--1000)\strut
\end{minipage} & \begin{minipage}[t]{0.20\columnwidth}\raggedright
Cubes, cube roots\strut
\end{minipage}\tabularnewline
\begin{minipage}[t]{0.28\columnwidth}\raggedright
S\strut
\end{minipage} & \begin{minipage}[t]{0.44\columnwidth}\raggedright
Sine scale, angle in degrees\strut
\end{minipage} & \begin{minipage}[t]{0.20\columnwidth}\raggedright
sin(θ) read against C/D\strut
\end{minipage}\tabularnewline
\begin{minipage}[t]{0.28\columnwidth}\raggedright
T\strut
\end{minipage} & \begin{minipage}[t]{0.44\columnwidth}\raggedright
Tangent scale, angle in degrees\strut
\end{minipage} & \begin{minipage}[t]{0.20\columnwidth}\raggedright
tan(θ) read against C/D\strut
\end{minipage}\tabularnewline
\bottomrule
\end{longtable}

Standard technique, if it's been a while: to multiply \emph{a × b}, set
the left (or right) index of C over \emph{a} on D, slide the cursor to
\emph{b} on C, read the result on D. Division reverses this. Squares:
set cursor on D, read A. Cubes: set cursor on D, read K. Sine/tangent:
set cursor to the angle on S (or T), read the corresponding value on C
(aligned back to D for a combined calculation).

\hypertarget{drill-1-multiplication-and-division-cd-scale}{%
\subsection{Drill 1 --- Multiplication and division (C/D
scale)}\label{drill-1-multiplication-and-division-cd-scale}}

Work these on the C/D scale only. Estimate the decimal placement
mentally \emph{before} you read the rule, then confirm.

\begin{enumerate}
\def\labelenumi{\arabic{enumi}.}
\tightlist
\item
  2.35 × 4.12
\item
  7.68 × 1.94
\item
  3.14 × 6.02
\item
  56.3 ÷ 8.15
\item
  128 ÷ 3.7
\item
  9.81 × 2.71828
\item
  0.0456 × 273
\item
  84.2 ÷ 0.0692
\item
  (4.5 × 6.3) ÷ 2.1 --- do this as one continuous setting, without
  resetting the rule between the multiply and the divide. That chaining
  is the actual skill; a real hand-computer never re-set the rule for an
  intermediate result if they could avoid it.
\end{enumerate}

\hypertarget{drill-2-squares-cubes-and-roots-ab-k-scale}{%
\subsection{Drill 2 --- Squares, cubes, and roots (A/B, K
scale)}\label{drill-2-squares-cubes-and-roots-ab-k-scale}}

\begin{enumerate}
\def\labelenumi{\arabic{enumi}.}
\tightlist
\item
  3.85²
\item
  12.7²
\item
  0.542²
\item
  2.15³
\item
  √56.3
\item
  √842
\item
  ∛47.5
\end{enumerate}

\hypertarget{drill-3-trig-scales-s-t}{%
\subsection{Drill 3 --- Trig scales (S,
T)}\label{drill-3-trig-scales-s-t}}

\begin{enumerate}
\def\labelenumi{\arabic{enumi}.}
\tightlist
\item
  sin(23.5°)
\item
  sin(67°)
\item
  tan(15°)
\item
  tan(52°)
\item
  45.0 × sin(30.0°) --- combined trig + multiply, one setting
\item
  68.3 ÷ tan(41.0°) --- combined trig + divide, one setting
\end{enumerate}

\begin{center}\rule{0.5\linewidth}{0.5pt}\end{center}

\hypertarget{worksheet}{%
\subsection{Worksheet}\label{worksheet}}

Fill in your slide rule reading and your decimal-placement estimate
\emph{before} checking the answer key.

\begin{longtable}[]{@{}lll@{}}
\toprule
Problem & Your estimate (order of magnitude) & Your slide rule
reading\tabularnewline
\midrule
\endhead
1.1 & &\tabularnewline
1.2 & &\tabularnewline
1.3 & &\tabularnewline
1.4 & &\tabularnewline
1.5 & &\tabularnewline
1.6 & &\tabularnewline
1.7 & &\tabularnewline
1.8 & &\tabularnewline
1.9 & &\tabularnewline
2.1 & &\tabularnewline
2.2 & &\tabularnewline
2.3 & &\tabularnewline
2.4 & &\tabularnewline
2.5 & &\tabularnewline
2.6 & &\tabularnewline
2.7 & &\tabularnewline
3.1 & &\tabularnewline
3.2 & &\tabularnewline
3.3 & &\tabularnewline
3.4 & &\tabularnewline
3.5 & &\tabularnewline
3.6 & &\tabularnewline
\bottomrule
\end{longtable}

\end{document}
